\chapter{Conclusion}
\label{ch:conclusion}

This thesis examined a single phenomenon from two sides. The
revaluation of AI-exposed equities since late 2022 posed a valuation
question with both a price dimension and a monetary one: whether the
prices of the S\&P~500 firms most engaged with artificial intelligence
have behaved in ways that ordinary fundamentals-driven pricing would
not imply, and whether those same firms are unusually sensitive to the
monetary-policy news that matters most for long-duration valuations.
Partitioning the index by pre-2023 10-K AI disclosure
\citep{anteSaggu2025}, it brought a distinct empirical design to bear
on each.

On the first question, the recursive right-tailed unit-root procedure
of \citet{phillipsShiYu2015a} finds explosive price dynamics that run
through the whole index but are ordered by AI exposure: the High AI
portfolio and its ratio to No AI firms reject the unit-root-with-drift
null at the $99\,\%$ level, the Moderate AI series only marginally, and
the sharpest signal lies in the Mag-7 concentration baskets. The
differential episodes against No AI firms fall mainly before ChatGPT,
with the clearest post-2022 signal in High-AI price levels and the
Mag-7 baskets, and the High AI cohort is entangled with a
longer-running pattern of technology-sector outperformance. The test
speaks to the shape of the price path, not to mispricing: it
identifies mildly explosive behaviour concentrated among AI-associated
firms, not a verdict that their prices are a bubble.

On the second question, the panel local projections of
\citet{jorda2005} on \citet{swanson2021} forward-guidance surprises
reveal a differential that is absent at impact and emerges only with
delay: relative to No AI firms in the same sector and on the same FOMC
date, a $+25$ basis-point hawkish surprise leaves High AI firms with
cumulative returns roughly six to eight percentage points lower over
the following two to three months, while the Moderate AI response is
indistinguishable from zero throughout. Because the
sector\,$\times$\,event fixed effects difference out the response
common to a sector at each event, this differential is identified
within sector, not across the technology-versus-rest divide; it
survives alternative shocks, outcomes, classifications and windows,
and a leverage-amplification check rules out financial leverage as its
source. What the design does not do is separate the valuation channels
through which the sensitivity operates.

Read together, the two exercises are where the thesis makes its
contribution. Each adds AI exposure to a literature that had not
included it --- to the cross-sectional bubble-detection evidence on
the current boom, previously confined to a single index and a fixed
mega-cap basket, and to the study of how firm characteristics shape
the transmission of monetary policy to equity prices, previously
focused on financial and structural traits. The more distinctive
result is the joint one: the same partition along which explosive
price dynamics are most pronounced is also the partition along which
the monetary-sensitivity differential is strongest. That
correspondence is reported as a pattern across two independent
designs, not as an identified mechanism.

Two questions remain open, and each points to a concrete next step. On
the pricing side, the explosive dynamics cannot yet be cleanly
separated from the technology-sector outperformance with which High AI
exposure is entangled; controlling for sector directly --- comparing
High AI firms with others in the same sector, or netting the AI
baskets against a technology benchmark rather than the broad No AI
group --- would isolate whatever explosive component is specific to AI,
much as the within-sector design already does on the monetary side. On
that monetary side, a present-value decomposition in the manner of
\citet{campbellAmmer1993} would begin to separate the discount-rate,
cash-flow and risk-premium channels that the present design leaves
bundled. Both exercises would, in turn, be sharpened by an exposure
measure that captures AI adoption rather than disclosure, such as the
workforce-based measure of \citet{eisfeldtSchubertZhangTaska2023},
against which the disclosure-based ranking used here could be
validated.

What the thesis establishes, then, is that AI exposure marks a
cross-section along which both explosive pricing and monetary
sensitivity have differed in the recent S\&P~500 --- robustly, on the
monetary side, and not through financial leverage. What it leaves to
further work is whether that cross-section is the footprint of
artificial intelligence itself or of the technology cycle that has
carried it: the question, in the end, with which it began.